\documentclass[11pt]{article}
\usepackage[margin=1in]{geometry}
\usepackage{amsmath,amssymb,amsthm}
\newtheorem{theorem}{Theorem}
\newtheorem{lemma}{Lemma}
\newcommand{\diag}{\operatorname{diag}}
\newcommand{\detm}{\operatorname{det}}
\newcommand{\tr}{\operatorname{tr}}
\title{Separated masses and reverse Schur pivots}
\author{Independent algebraic audit for the bounded-degree construction}
\date{November 15, 2026}
\begin{document}
\maketitle

\begin{theorem}[General separated-mass fractional limit]\label{thm:main}
Let $K\in\mathbb R^{d\times d}$ be positive definite, let
$K'=K+vv^T$, and fix $0<\alpha<1$. Put
\[
 D_\delta=\diag(\delta^{(d-1)/2},\delta^{(d-2)/2},\ldots,1),\quad
 A_\delta=D_\delta K D_\delta,\quad
 B_\delta=D_\delta K'D_\delta,\qquad\delta\downarrow0.
\]
For $T_i=\{i,\ldots,d\}$ define
\[
 \tau_i=v_{T_i}^T K_{T_i,T_i}^{-1}v_{T_i},\quad
 \tau_{d+1}=0,\qquad
 \kappa_i=\frac{\det K_{T_i,T_i}}{\det K_{T_{i+1},T_{i+1}}},
\]
where an empty determinant equals one. Define $\kappa_i'$ in the same
way using $K'$. Then, in a suitably ordered orthonormal eigenbasis of
$A_\delta$,
\[
 A_\delta^{-\alpha/2}(B_\delta^\alpha-A_\delta^\alpha)
                 A_\delta^{-\alpha/2}
 \longrightarrow
 \diag\left(\left(\frac{\kappa_i'}{\kappa_i}\right)^\alpha-1\right)_{i=1}^d
\]
in operator norm. Equivalently, the $i$th limiting relative factor is
\[
 \frac{\kappa_i'}{\kappa_i}
          =\frac{1+\tau_i}{1+\tau_{i+1}}.
\]
The result holds also for convergent families $K_\delta\to K\succ0$
and $v_\delta\to v$, with the same limiting formula.
\end{theorem}

\paragraph{Interpretation.}
The mass matrix is $M_\delta=D_\delta^{-2}$, so its masses obey
$m_1\gg m_2\gg\cdots\gg m_d$. The successive slow modes follow reverse
Schur elimination. No tridiagonal, tree or special update-coordinate
assumption is required.

\begin{lemma}[Eigenvalue scales]\label{lem:scales}
Writing the eigenvalues of $A_\delta$ in increasing order as
$0<\lambda_1(\delta)\le\cdots\le\lambda_d(\delta)$, one has
\[
 \lambda_i(\delta)\sim\kappa_i\delta^{d-i}.
\]
\end{lemma}
\begin{proof}
Let $e_k$ be the $k$th elementary symmetric polynomial of the
eigenvalues, and set $e_0=1$. The principal-minor formula gives
\[
 e_k(A_\delta)
   =\sum_{|S|=k}\det K_{S,S}\,\delta^{\sum_{j\in S}(d-j)}
   \sim\det K_{T_{d-k+1},T_{d-k+1}}\,\delta^{k(k-1)/2}.
\]
The dominating index set is unique. Let
$P_k=\lambda_{d-k+1}\cdots\lambda_d$.
Since $P_k\le e_k\le \binom dk P_k$, for $1\le k<d$,
\[
 \frac{\lambda_{d-k}}{\lambda_{d-k+1}}
  =\frac{P_{k+1}P_{k-1}}{P_k^2}
  \le \binom dk^2\frac{e_{k+1}e_{k-1}}{e_k^2}
  =O(\delta).
\]
Thus adjacent scales separate. Every product contributing to $e_k$
other than $P_k$ is $o(P_k)$, so $e_k\sim P_k$.
Taking the ratio of $e_{d-i+1}$ and $e_{d-i}$ gives the claimed
constant as well as the exponent.
\end{proof}

\begin{lemma}[Leverage weights on separated scales]\label{lem:weights}
Let $u_\delta=D_\delta v$, choose orthonormal eigenvectors $\phi_i$ of
$A_\delta$ corresponding to the ordering in Lemma~\ref{lem:scales},
and set
\[
 z_i(\delta)=\frac{|\langle u_\delta,\phi_i\rangle|^2}{\lambda_i(\delta)}.
\]
Then $z_i(\delta)\to\tau_i-\tau_{i+1}$.
\end{lemma}
\begin{proof}
The total leverage is exactly fixed:
$\sum_i z_i=v^TK^{-1}v=\tau_1$. Thus the vector of nonnegative weights
lies in a compact simplex. Consider any convergent subsequence,
with limits $z_i^*$. Define the scalar resolvent response
\[
 g_\delta(t)=u_\delta^T(A_\delta+tI)^{-1}u_\delta
           =v^T(K+tD_\delta^{-2})^{-1}v.
\]
At $t=\delta^{d-i}s$, where $s>0$ is fixed, the extra diagonal terms
at indices $j<i$ diverge, the term at $i$ equals $s$, and the terms
at $j>i$ vanish. Block elimination therefore gives
\[
 g_\delta(\delta^{d-i}s)
 \longrightarrow
 v_{T_i}^T(K_{T_i,T_i}+s e_ie_i^T)^{-1}v_{T_i}.
\]
Writing $K_{T_i,T_i}$ in the block decomposition $\{i\}\cup T_{i+1}$,
its scalar Schur complement is $\kappa_i$. If
\[
 h_i=v_i-K_{i,T_{i+1}}K_{T_{i+1},T_{i+1}}^{-1}v_{T_{i+1}},
\]
the limit is
\[
 \tau_{i+1}+\frac{h_i^2}{\kappa_i+s},\qquad
 \tau_i-\tau_{i+1}=\frac{h_i^2}{\kappa_i}.
\]
On the other hand, the spectral formula and Lemma~\ref{lem:scales}
give along the chosen subsequence
\[
 g_\delta(\delta^{d-i}s)
   =\sum_j\frac{\lambda_j z_j}{\lambda_j+\delta^{d-i}s}
   \longrightarrow\sum_{j>i}z_j^*
                      +\frac{\kappa_i z_i^*}{\kappa_i+s}.
\]
Equality for every $s>0$ identifies $z_i^*=h_i^2/\kappa_i$.
The limit is unique for every subsequence, proving the assertion.
\end{proof}

\begin{proof}[Proof of Theorem~\ref{thm:main}]
Use the eigenbasis from Lemma~\ref{lem:weights}, adjusting signs of
the eigenvectors so the coordinates of $u_\delta$ are nonnegative.
Set $c_\alpha=\sin(\pi\alpha)/\pi$. The exact fractional resolvent
formula gives the $i$th relative diagonal entry, after $t=\lambda_i s$,
as
\[
 C_{ii}=c_\alpha z_i\int_0^\infty\frac{s^\alpha}{(1+s)^2}
 \left(1+\sum_j\frac{\lambda_j z_j}{\lambda_j+\lambda_i s}\right)^{-1}ds.
\]
The scalar resolvent term in parentheses tends to
\[
 1+\sum_{j>i}(\tau_j-\tau_{j+1})
        +\frac{\tau_i-\tau_{i+1}}{1+s}
 =1+\tau_{i+1}+\frac{\tau_i-\tau_{i+1}}{1+s}.
\]
The weights are uniformly bounded and $s^\alpha/(1+s)^2$ is integrable.
Dominated convergence and the scalar fractional formula therefore give
\[
 C_{ii}\longrightarrow
 \left(\frac{1+\tau_i}{1+\tau_{i+1}}\right)^\alpha-1.
\]
Omitting the positive Sherman--Morrison denominator gives, for $i<j$,
\[
 |C_{ij}|\le\sqrt{z_iz_j}\,x^{(1-\alpha)/2}
                     \frac{1-x^\alpha}{1-x},\qquad x=\lambda_i/\lambda_j.
\]
The scale lemma makes $x\to0$, so every off-diagonal entry tends to
zero. Fixed dimension upgrades entrywise convergence to norm convergence.
Finally, the determinant lemma on each trailing principal submatrix gives
\[
 \frac{\kappa_i'}{\kappa_i}
 =\frac{\det K'_{T_i}/\det K_{T_i}}
        {\det K'_{T_{i+1}}/\det K_{T_{i+1}}}
 =\frac{1+\tau_i}{1+\tau_{i+1}}.
\]
All arguments use finite sums, matrix inversion with a positive
definite limiting matrix, and uniform integrable domination. They
therefore apply unchanged when $K_\delta\to K\succ0$ and
$v_\delta\to v$.
\end{proof}

\paragraph{How the fanout construction uses the lemma.}
Suppose the Kron-reduced terminal pair satisfies $K'=K+vv^T$ and,
in the reservoir elimination order, one proves
\[
 \frac{\det K'_{T_i}/\det K_{T_i}}
      {\det K'_{T_{i+1}}/\det K_{T_{i+1}}}
       \ge1+\frac{c}{m}
       \qquad(1\le i\le m)
\]
for an absolute $c>0$. Taking the reservoir masses sufficiently
separated then yields $m$ relative fractional correction eigenvalues
at least
\[
 \frac12\left[\left(1+\frac cm\right)^\alpha-1\right]
       \ge\frac{\alpha c}{2(1+c)m}.
\]
The last inequality follows by integrating
$\alpha(1+x)^{\alpha-1}\ge\alpha/(1+c)$ over $0\le x\le c/m$.
Thus each finite $m$ can be realized with a finite mass hierarchy.
The graph-specific remaining obligation is the reservoir reduction
from ordinary unit masses; it is distinct from the algebraic lemma
proved here.
\end{document}
